\section{Vocabulaire des opérations}
\begin{vocab}[(Copier AVEC DES COULEURS)]
\includegraphics[scale=0.4]{ASSETS/vocabulaire}
\end{vocab}


\begin{vocab}[(Pas obligatoire de copier si l'exemple ci-dessus vous suffit)]
\begin{itemize}
    \item Le résultat d'une \hlb[RougeClair]{addition}
        (\hlb[VertClair]{soustraction}) s'appelle une \answer[RougeClair]{\bb{somme}}
        (\answer[VertClair]{\bb{différence}}) et les nombres utilisés s'appellent les
        \answer[GrisClair]{termes}.
    \item Le résultat d'une \hlb[BleuClair]{multiplication} s'appelle un
        \answer[BleuClair]{\bb{produit}} et les nombres utilisés s'appellent les
        \answer[BleuClair]{\bb{facteurs}}.
    \item Le résultat d'une \hlb[JauneClair]{division} s'appelle un
        \answer[JauneClair]{\bb{quotient}} et les nombres utilisés s'appellent le
        \answer[JauneClair]{\bb{dividende}} et le \answer[JauneClair]{\bb{diviseur}}.
        (\answer[JauneClair]{$$\text{dividende} \div \text{diviseur} = \text{quotient}$$})
\end{itemize}
\end{vocab}

\begin{remarquebox}
    Selon la dernière opération effectuée, on dit qu'une expression est une
    \bb{somme}, un \bb{produit}, une \bb{différence} ou un \bb{quotient}.
    \newline
    \newline
    Explication en vidéo :
    \href{https://app.videas.fr/ae1ff766-dd4e-4de1-9efe-f8ae2b7992e7/}{ici}
\end{remarquebox}

\begin{exemplebox}
    $5+2\times3$ \hspace{5pt}est \answer{une somme} car la dernière opération effectuée est
    \answer{une addition}. 
    La traduction en phrase de cette expression est
    \answer{La somme de 5 et du produit de 2 par 3}.

    $(2+3)\times5$ est \answer{un produit} car la dernière opération effectuée
    est \answer{une multiplication}. La traduction en phrase de cette expression
    est \answer{\it{Le produit de la somme de 2 et 3 par 5}}
\end{exemplebox}

